% !TeX spellcheck = en_US
\chapter{Experiments and Results}%
\label{ch:experiments}%

This chapter describes the experimental realization of the methods introduced in \cref{ch:methods}.
\Cref{sec:experimental_setup} presents the experimental setup, covering teacher policy training (\cref{sec:teacher_training_setup}), student policy distillation (\cref{sec:student_implementation}), and the evaluation metrics (\cref{sec:eval_protocol}).
\Cref{sec:ablations} describes the ablation studies and their rationale.
\Cref{sec:results} reports the results for both teacher training and distillation.

%% =========================================================================
\section{Experimental Setup}%
\label{sec:experimental_setup}%

All experiments are conducted in Isaac Lab 2.1.0~\cite{Mittal2023} (Isaac Sim 5.1.0) on a single NVIDIA RTX 4090 GPU.
The simulation runs at 120\,Hz with a policy control frequency of 60\,Hz.
The robotic platform consists of a \ac{FR3} 7-\acs{DoF} arm equipped with a 16-\acs{DoF} Agile Hand (Agile Robots SE), forming a 23-\acs{DoF} arm--hand system controlled via joint position targets with simulated PD controllers (\cref{fig:hardware_platform}).
The arm and hand are combined into a single articulated \ac{USD} asset with convex-hull collision meshes.
The observation, action, and reward design are defined in \cref{sec:asymmetric_ac,sec:reward_structure}.

\begin{figure}[htb]
    \centering
    \begin{subfigure}[b]{0.45\textwidth}
        \centering
        \includegraphics[width=\textwidth]{setup/Agile_Hand_photo}
        \caption{The physical Agile Hand.}
        \label{fig:agile_hand_photo}
    \end{subfigure}
    \hfill
    \begin{subfigure}[b]{0.45\textwidth}
        \centering
        \includegraphics[width=\textwidth]{setup/Setup_isaacsim}
        \caption{FR3 + Agile Hand in Isaac Sim.}
        \label{fig:setup_isaacsim}
    \end{subfigure}
    \caption{Hardware platform: (a)~the 16-\acs{DoF} Agile Hand manufactured by Agile Robots SE, and (b)~the combined FR3 arm--hand system as rendered in the Isaac Sim simulation environment.}
    \label{fig:hardware_platform}
\end{figure}

\subsection{Teacher Policy Training}%
\label{sec:teacher_training_setup}%

The teacher is trained with \ac{PPO} on 1{,}024 parallel environments with \ac{ADR} following the formulation in \cref{sec:adr_formulation}.
We evaluate on 13 objects from the VisDex dataset (\cref{fig:training_objects}), spanning compact graspable shapes (teddy bear, closed fist) to thin and irregular objects (plane, toy bagger).
At episode reset, each environment is assigned one object uniformly at random, and the object's mass, friction, and scale are randomized within the \ac{ADR} ranges.
Each episode lasts up to 10\,s (600 steps) and terminates early if any unsafe condition from \cref{tab:terminations} is triggered.
These termination conditions are reused as unsafe-event signals for SafeDAgger's intervention logic during distillation.
An episode is considered successful if the object is lifted above the table and held for at least 0.5\,s (30 steps).
When the rolling success rate across all environments exceeds the \ac{ADR} threshold $\tau_{\text{ADR}} = 0.4$ (40\,\%), the randomization ranges are widened by one increment, progressively increasing the difficulty of the training distribution.
Training hyperparameters are listed in \cref{app:ppo_hyperparams}.

\begin{figure}[htb]
    \centering
    \begin{subfigure}[b]{0.24\textwidth}
        \centering
        \includegraphics[width=\textwidth]{trained_objects/mario_view_1}
        \caption{Mario}
    \end{subfigure}\hfill
    \begin{subfigure}[b]{0.24\textwidth}
        \centering
        \includegraphics[width=\textwidth]{trained_objects/milk_pot_view_1}
        \caption{Milk pot}
    \end{subfigure}\hfill
    \begin{subfigure}[b]{0.24\textwidth}
        \centering
        \includegraphics[width=\textwidth]{trained_objects/turtle_candle_tolder_view_1}
        \caption{Turtle candle holder}
    \end{subfigure}\hfill
    \begin{subfigure}[b]{0.24\textwidth}
        \centering
        \includegraphics[width=\textwidth]{trained_objects/toy_bagger_view_1}
        \caption{Toy bagger}
    \end{subfigure}

    \vspace{0.5em}
    \begin{subfigure}[b]{0.24\textwidth}
        \centering
        \includegraphics[width=\textwidth]{trained_objects/teddy_bear_view_1}
        \caption{Teddy bear}
    \end{subfigure}\hfill
    \begin{subfigure}[b]{0.24\textwidth}
        \centering
        \includegraphics[width=\textwidth]{trained_objects/closed_fist_view_1}
        \caption{Closed fist}
    \end{subfigure}\hfill
    \begin{subfigure}[b]{0.24\textwidth}
        \centering
        \includegraphics[width=\textwidth]{trained_objects/basketball_shoe_view_1}
        \caption{Basketball shoe}
    \end{subfigure}\hfill
    \begin{subfigure}[b]{0.24\textwidth}
        \centering
        \includegraphics[width=\textwidth]{trained_objects/chickenhead_in_car_view_1}
        \caption{Chicken head in car}
    \end{subfigure}

    \vspace{0.5em}
    \begin{subfigure}[b]{0.24\textwidth}
        \centering
        \includegraphics[width=\textwidth]{trained_objects/toy_train_view_1}
        \caption{Train}
    \end{subfigure}\hfill
    \begin{subfigure}[b]{0.24\textwidth}
        \centering
        \includegraphics[width=\textwidth]{trained_objects/homer_view_1}
        \caption{Homer}
    \end{subfigure}\hfill
    \begin{subfigure}[b]{0.24\textwidth}
        \centering
        \includegraphics[width=\textwidth]{trained_objects/toy_cow_view_1}
        \caption{Toy cow}
    \end{subfigure}\hfill
    \begin{subfigure}[b]{0.24\textwidth}
        \centering
        \includegraphics[width=\textwidth]{trained_objects/plane_view_1}
        \caption{Plane}
    \end{subfigure}\hfill
    \begin{subfigure}[b]{0.24\textwidth}
        \centering
        \includegraphics[width=\textwidth]{trained_objects/elefant_toy_view_1}
        \caption{Elephant toy}
    \end{subfigure}
    \caption{The 13 training objects from the VisDex dataset.}
    \label{fig:training_objects}
\end{figure}

\begin{table}[htb]
    \centering
    \caption{Episode termination conditions. All early terminations mark the state as unsafe; penalty-eligible conditions additionally trigger $r_{\text{term}}$ when they occur before any object contact.}
    \label{tab:terminations}
    \small
    \begin{tabular}{lp{7.5cm}c}
        \toprule
        \textbf{Condition} & \textbf{Trigger} & \textbf{Penalty} \\
        \midrule
        \texttt{object\_escape\_xy}    & Object leaves the $\pm0.25\,\mathrm{m}\times\pm0.4\,\mathrm{m}$ workspace rectangle & \checkmark \\
        \texttt{object\_too\_low}      & Object falls below the table surface ($z < 0.2\,\mathrm{m}$) & \checkmark \\
        \texttt{hand\_too\_far}        & Hand workspace body leaves a $0.35\,\mathrm{m}$-inflated workspace box & \checkmark \\
        \texttt{hand\_too\_close}      & Any fingertip or palm comes within $1\,\mathrm{cm}$ of the table surface & -- \\
        \texttt{palm\_flip}            & Palm direction deviates more than $108^{\circ}$ from the target down-direction & \checkmark \\
        \texttt{arm\_table\_contact}   & Any arm link contacts the table & -- \\
        \texttt{hand\_vel\_explosion}  & Any finger joint velocity exceeds $50\,\mathrm{rad/s}$ & -- \\
        \texttt{robot\_unstable}       & NaN or Inf detected in joint positions or velocities & -- \\
        \midrule
        \texttt{timeout}               & Episode reaches $T_{\max} = 10\,\mathrm{s}$ (600 steps at $60\,\mathrm{Hz}$) & -- \\
        \bottomrule
    \end{tabular}
\end{table}

\subsection{Student Policy Distillation}%
\label{sec:student_implementation}%

The student replaces the teacher's privileged object state with stereo RGB images ($320 \times 240$ at 30\,fps) while retaining the proprioceptive inputs (joint positions/velocities, hand pose, goal position, previous action).
The student does \emph{not} receive the object position, orientation, scale, or contact information; these must be inferred from the visual input.
The visual encoder, cross-attention fusion, and LSTM policy head are described in \cref{sec:rgb_distillation}.

Distillation follows the SafeDAgger procedure from \cref{alg:safedagger}, using Adam with a learning rate of $1 \times 10^{-4}$ and the combined loss from \cref{eq:total_loss}.
The teacher intervenes whenever the per-environment L2 action disagreement exceeds the safety threshold $\delta$ (\cref{sec:ablation_threshold}).
Training runs for 100{,}000 iterations across 24 parallel environments.
The object set is reduced from the teacher's 13 VisDex objects to an 8-object subset (basketball shoe, closed fist, elephant toy, mario, milk pot, teddy bear, toy bagger, turtle candle holder) to ensure at least 3 environments per object, which produces more stable per-object gradients during online distillation.

%% =========================================================================
\subsection{Evaluation Metrics}%
\label{sec:eval_protocol}%

Three metrics are used to evaluate both teacher and student policies.
Let $\tau \in \{1, \ldots, N\}$ index evaluation episodes and $t \in \{0, \ldots, T_\tau\}$ index time steps within an episode.
Unless otherwise stated, evaluation episodes use a horizon of 10\,s and each configuration is evaluated over $N = 480$--$640$ episodes (60--80 per object).

\paragraph{Lift success (grasp with hold).}
Let $h_t^\tau$ denote the object height above the table at time step $t$ of episode $\tau$, and let $H$ be a lift threshold.
An episode is considered successful if the object remains above $H$ for at least $N_{\text{hold}}$ consecutive steps:
%
\begin{equation}\label{eq:lift_success}
    \mathbb{I}^{\tau}_{\text{lift}} = \mathbb{I}\!\left[\exists\, t : \forall\, k \in \{0, \ldots, N_{\text{hold}}-1\},\; h_{t+k}^{\tau} > H\right].
\end{equation}
%
In this work, $H$ corresponds to the table surface and $N_{\text{hold}} = 30$ (0.5\,s at 60\,Hz).
The hold requirement prevents counting accidental upward launches as successful grasps.
The \ac{GSR} is the primary task metric:
%
\begin{equation}\label{eq:gsr}
    \text{GSR} = \frac{1}{N} \sum_{\tau=1}^{N} \mathbb{I}^{\tau}_{\text{lift}}.
\end{equation}

\paragraph{Unsafe episode rate.}
An episode may terminate early due to one of four hardware-relevant failure modes:
\begin{itemize}
    \item \textbf{Harmful collision} ($m = \text{coll}$): any fingertip or palm body comes within 1\,cm of the table surface.
    \item \textbf{Object out of bounds} ($m = \text{obj}$): the object leaves the workspace rectangle or falls below 0.2\,m height.
    \item \textbf{Palm flipped} ($m = \text{flip}$): the palm direction deviates more than $108^{\circ}$ from the downward-facing target direction.
    \item \textbf{Hand too far} ($m = \text{hf}$): the hand base link leaves a 0.35\,m-inflated bounding box around the table workspace.
\end{itemize}
Let $\mathcal{M} = \{\text{coll}, \text{obj}, \text{flip}, \text{hf}\}$ denote this set of termination modes, and let $\mathbb{I}^{\tau}_{m} \in \{0,1\}$ indicate whether mode $m \in \mathcal{M}$ occurs in episode $\tau$.
An episode is unsafe if any mode is triggered:
%
\begin{equation}\label{eq:unsafe}
    \mathbb{I}^{\tau}_{\text{unsafe}} = \mathbb{I}\!\left[\bigvee_{m \in \mathcal{M}} \mathbb{I}^{\tau}_{m} = 1\right].
\end{equation}
%
The \ac{UER} quantifies training-time safety:
%
\begin{equation}\label{eq:uer}
    \text{UER} = \frac{1}{N} \sum_{\tau=1}^{N} \mathbb{I}^{\tau}_{\text{unsafe}}.
\end{equation}
%
This definition excludes \texttt{physics\_instability} events (NaN joint states, velocity explosions), which are simulation-only artifacts that cannot occur on physical hardware.
All \ac{UER} values reported in this thesis refer to this hardware-relevant definition unless stated otherwise.
Full results including physics instability are provided in \cref{app:full_results}.

\paragraph{Failure-mode decomposition.}
Each episode is assigned a single terminal mode $m^\tau \in \mathcal{M} \cup \{\text{timeout}, \text{success}\}$.
When multiple unsafe conditions trigger in the same time step, the first condition (by a fixed priority order) is recorded.
Per-mode rates are reported as:
%
\begin{equation}\label{eq:failure_mode}
    \text{FailRate}(m) = \frac{1}{N} \sum_{\tau=1}^{N} \mathbb{I}[m^\tau = m], \quad m \in \mathcal{M}.
\end{equation}
%
This decomposition identifies which failure categories each distillation method is most prone to.

%% =========================================================================
\section{Ablation Studies}%
\label{sec:ablations}%

The following ablation studies are conducted to justify the pipeline design.
The trained teacher policy is distilled into a vision-based student using three distillation variants: \ac{BC}, DAgger, and SafeDAgger.
All three use the same teacher, student architecture, loss function, and training budget; only the data-collection strategy differs.
The variants are evaluated in terms of lift-success growth, unsafe-episode rate, and intervention rate (defined as the fraction of time steps in which teacher takeover occurs, $\beta = \frac{1}{T}\sum_{t=0}^{T-1}\mathbb{I}[g_t = 1]$, averaged over evaluation episodes) during training.
Since the \ac{UER} is the primary safety metric, not only the average rate across all test objects is reported but also a per-object, failure-mode decomposition to assess each method under different failure conditions.
Additionally, standalone student evaluations provide the final task performance after training.

\subsection{Safety Threshold Sensitivity}%
\label{sec:ablation_threshold}%

SafeDAgger's behavior depends on the safety threshold $\delta$, which Zhang and Cho~\cite{Zhang2016} selected as a single value based on the error distribution of an initial supervised policy without systematic analysis.
Following the principle from the original paper, $\delta$ should be calibrated so that approximately 20--30\,\% of environments trigger teacher intervention at the start of distillation, allowing the intervention rate to decay naturally as the student improves.

The calibration proceeds by analyzing the distribution of per-environment L2 action disagreement $e_t = \|\hat{\pi}_\theta(o_t^{\text{student}}) - \pi^*(o_t^{\text{teacher}})\|_2$ during early training and selecting $\delta$ at a percentile that produces the target intervention rate.
This work sweeps $\delta$ across the values 0.5, 2.0, and 3.0, with $\delta \to \infty$ recovering vanilla DAgger as the upper bound.
The results of this sweep are reported in \cref{sec:threshold_results}.

\subsection{Distillation Method Comparison}%
\label{sec:ablation_method}%

The central experiment compares the three distillation methods using the threshold identified from the sensitivity sweep (\cref{sec:ablation_threshold}).
\ac{BC} trains on a fixed offline dataset of teacher rollouts, DAgger iteratively collects data under the student's own policy, and SafeDAgger adds safety-aware gating (\cref{alg:safedagger}).

%% =========================================================================
\section{Results, Graphs and Tables}%
\label{sec:results}%

\subsection{Teacher Training Results}%
\label{sec:teacher_results}%

The teacher policy was trained on 13 objects from the VisDex dataset using 1{,}024 parallel environments with a sim-to-real-aware \ac{ADR} curriculum.
Training reached ADR level~14 after approximately 20{,}000 epochs and continued for 33{,}000 epochs total without further ADR progression.
The resulting teacher achieves 85.8\,\% \ac{GSR} and 12.1\,\% \ac{UER} at nominal (ADR-0) physics, demonstrating reliable multi-object grasping with the hardware-realistic actuator constraints.
\Cref{tab:teacher_eval} reports the full evaluation breakdown.
This teacher is used for all distillation experiments reported below.

\begin{table}[htb]
    \centering
    \caption{Teacher policy per-object evaluation (200 episodes per object, ADR-0 physics). Two objects (chicken head in car, train) failed to learn grasping.}
    \label{tab:teacher_eval}
    \small
    \begin{tabular}{lrrrrr}
        \toprule
        \textbf{Object} & \textbf{Lift (\%)} & \textbf{UER (\%)} & \textbf{Coll. (\%)} & \textbf{OOB (\%)} & \textbf{Palm (\%)} \\
        \midrule
        basketball shoe       &  99.0 & 28.0 &  7.0 & 21.0 &  0.0 \\
        closed fist           & 100.0 &  9.0 &  0.0 &  9.0 &  0.0 \\
        elephant toy          & 100.0 &  3.0 &  2.0 &  0.5 &  0.5 \\
        homer                 &  99.0 & 12.5 &  2.0 & 10.5 &  0.0 \\
        mario                 &  98.0 &  2.0 &  0.0 &  1.0 &  0.0 \\
        milk pot              &  96.5 &  5.5 &  0.5 &  5.0 &  0.0 \\
        plane                 &  83.5 & 12.0 &  1.5 &  2.5 &  8.0 \\
        teddy bear            &  99.5 & 15.0 &  0.5 & 14.5 &  0.0 \\
        toy bagger            &  99.0 & 10.5 &  2.0 &  7.5 &  1.0 \\
        toy cow               & 100.0 & 12.5 &  0.0 & 12.0 &  0.5 \\
        turtle candle holder  & 100.0 & 24.5 &  0.0 & 24.0 &  0.0 \\
        \midrule
        chicken head in car   &   0.0 &  0.0 &  0.0 &  0.0 &  0.0 \\
        train                 &   0.0 & 11.0 &  2.0 &  7.5 &  0.5 \\
        \midrule
        \textbf{Average (11 lifted)} & \textbf{97.7} & \textbf{12.2} & \textbf{1.4} & \textbf{9.8} & \textbf{0.9} \\
        \textbf{Average (all 13)}    & \textbf{82.7} & \textbf{11.2} & \textbf{1.3} & \textbf{8.8} & \textbf{0.8} \\
        \bottomrule
    \end{tabular}
\end{table}

%% =========================================================================
\subsection{Safety Threshold Results}%
\label{sec:threshold_results}%

\Cref{tab:threshold_sweep} summarizes the effect of the safety threshold $\delta$ on SafeDAgger distillation performance.
An initial threshold of $\delta = 0.5$ (scaled from the original paper's steering-angle domain to the 23-dimensional action space) produced an intervention rate of $\sim$95\,\%, effectively reducing SafeDAgger to behavioral cloning and limiting \ac{GSR} to 61.3\,\%.
At $\delta = 3.0$, the intervention rate dropped to $\sim$5\,\%, making SafeDAgger behave essentially like DAgger, with a \ac{GSR} of only 27.9\,\%.
Analysis of the L2 distribution at early and late training (80th percentile: 2.6 early, 1.7 late) confirmed that $\delta = 2.0$ produces $\sim$30\,\% initial intervention decaying to $\sim$5\,\% at convergence, matching the target regime from the original SafeDAgger paper.
This value achieves the highest \ac{GSR} (72.7\,\%) while maintaining the lowest \ac{UER} (30.2\,\%) and is used for all subsequent distillation experiments.

\begin{table}[htb]
    \centering
    \caption{Effect of safety threshold $\delta$ on SafeDAgger distillation performance (100k iterations, 24 environments, 480 standalone evaluation episodes).}
    \label{tab:threshold_sweep}
    \begin{tabular}{lccc}
        \toprule
        \textbf{$\delta$} & \textbf{Method} & \textbf{\acs{GSR} (\%)} & \textbf{\acs{UER} (\%)} \\
        \midrule
        0.5       & SafeDAgger & 61.3 & 34.8 \\
        2.0       & SafeDAgger & 72.7 & 30.2 \\
        3.0       & SafeDAgger & 27.9 & 26.9 \\
        $\infty$  & DAgger     & 71.7 & 32.7 \\
        \bottomrule
    \end{tabular}
\end{table}

\Cref{fig:threshold_sweep} shows the intervention rate $\beta_{\text{safe}}$ over training for each threshold.
At $\delta = 0.5$, the teacher intervenes in $\sim$70--80\,\% of environments throughout training, preventing the student from practicing autonomously.
At $\delta = 2.0$, the intervention rate decays from $\sim$30\,\% to $\sim$10\,\% within 20k iterations, indicating rapid student autonomy.
At $\delta = 3.0$, intervention is minimal from the start ($<$10\,\%), offering little safety benefit over vanilla DAgger.

\begin{figure}[htbp]
    \centering
    \includegraphics[width=0.75\textwidth]{plots/safedagger_threshold/intervention_rate_threshold_sweep}
    \caption{SafeDAgger intervention rate $\beta_{\text{safe}}$ over training iterations for different safety thresholds $\delta \in \{0.5, 2.0, 3.0\}$. Lower thresholds cause more aggressive teacher intervention. At $\delta = 2.0$, the rate decays rapidly, balancing safety with student autonomy.}
    \label{fig:threshold_sweep}
\end{figure}

%% =========================================================================
\subsection{Distillation Results}%
\label{sec:distillation_results}%

All distillation experiments share the same teacher (85.8\,\% \ac{GSR}, 12.1\,\% \ac{UER}), the same student architecture (ResNet-18 + cross-attention + LSTM-512), and the same loss (\cref{eq:total_loss}).
The primary three-way comparison uses 100{,}000 distillation iterations on the 8-object VisDex subset.

\subsubsection{Distillation Method Comparison}%
\label{sec:dagger_results}%

\Cref{tab:distillation_comparison} reports the evaluation metrics for BC, DAgger, and SafeDAgger ($\delta = 2.0$).

\begin{table}[htb]
    \centering
    \caption{BC vs.\ DAgger vs.\ SafeDAgger standalone evaluation (480 episodes).}
    \label{tab:distillation_comparison}
    \begin{tabular}{lcccc}
        \toprule
        \textbf{Metric} & \textbf{BC} & \textbf{DAgger} & \textbf{SafeDAgger} & \textbf{Teacher} \\
        \midrule
        \acs{GSR} (\%) & 87.8 & 71.7 & 72.7 & 85.8 \\
        \acs{UER} (\%) & 41.8 & 32.7 & 30.2 & 12.1 \\
        \quad Harmful collision (\%) & 19.3 & 12.5 & 16.2 & 1.7 \\
        \quad Object out of bounds (\%) & 21.1 & 19.2 & 11.5 & 9.4 \\
        \quad Palm flipped (\%) & 1.3 & 1.1 & 2.5 & 0.6 \\
        \quad Hand too far (\%) & 0.0 & 0.0 & 0.0 & 0.0 \\
        \bottomrule
    \end{tabular}
\end{table}

BC achieves the highest \ac{GSR} (87.8\,\%), matching the teacher within 2 percentage points.
This is because BC trains exclusively on the teacher's state distribution, where the action-prediction problem is well-posed: the student learns to reproduce the teacher's actions with high fidelity for states the teacher actually visits.
DAgger and SafeDAgger, by contrast, let the student act in the environment during training, exposing it to its own imperfect trajectories and forcing it to learn from a harder, shifting distribution.
Because the lift success metric requires the object to be held above the threshold for at least 0.5\,s (30 consecutive steps), BC's high \ac{GSR} reflects genuinely sustained grasps, not brief accidental lifts.
However, BC's \ac{UER} is substantially higher (41.8\,\% vs.\ 30.2\,\% for SafeDAgger and 32.7\,\% for DAgger): the student learns to grasp well but cannot maintain safe behavior when it inevitably drifts from the teacher's state distribution.
Since the student never encounters its own out-of-distribution states during training, it does not learn recovery behaviors, resulting in a classic covariate-shift failure: high task accuracy paired with poor robustness at deployment.

A checkpoint sweep (\cref{tab:bc_overfitting}) confirms that BC overfits: lift success peaks at 95.6\,\% after 50k iterations, then degrades to 85.6\,\% at 100k.
The \ac{UER} remains consistently high (35--49\,\%) across all checkpoints, confirming that BC cannot learn safety behavior regardless of training duration.

\begin{table}[htb]
    \centering
    \caption{BC checkpoint sweep (640 episodes per checkpoint, 32 environments). Lift success peaks at 50k iterations and degrades thereafter, while \acs{UER} remains consistently high.}
    \label{tab:bc_overfitting}
    \begin{tabular}{rcccccc}
        \toprule
        \textbf{Iter} & \textbf{Lift (\%)} & \textbf{UER (\%)} & \textbf{Coll. (\%)} & \textbf{OOB (\%)} & \textbf{Palm (\%)} \\
        \midrule
        10k  & 71.9 & 49.4 & 13.4 & 12.5 & 23.4 \\
        25k  & 86.9 & 47.2 & 13.9 & 23.4 &  9.8 \\
        \textbf{50k}  & \textbf{95.6} & 39.1 & 19.4 & 16.6 &  3.1 \\
        75k  & 95.2 & 35.8 & 14.4 & 14.8 &  6.6 \\
        100k & 85.6 & 40.1 & 20.6 & 18.0 &  1.6 \\
        \bottomrule
    \end{tabular}
\end{table}

SafeDAgger and DAgger reach comparable \ac{GSR} (72.7\,\% vs.\ 71.7\,\%), recovering roughly 85\,\% of the teacher ceiling.
SafeDAgger maintains a lower \ac{UER} throughout training ($\sim$15--20\,\% vs.\ $\sim$30--40\,\% for DAgger, \cref{fig:run9_unsafe_rate}), with the largest difference in the object out of bounds category (11.5\,\% vs.\ 19.2\,\%).
\Cref{fig:run9_training_curves} shows the training curves and intervention rate for all three methods.

\begin{figure}[htbp]
    \centering
    \begin{subfigure}[b]{0.48\textwidth}
        \centering
        \includegraphics[width=\textwidth]{plots/run9/lift_success_normed}
        \caption{Normalized lift success (relative to teacher).}
        \label{fig:run9_lift_success}
    \end{subfigure}
    \hfill
    \begin{subfigure}[b]{0.48\textwidth}
        \centering
        \includegraphics[width=\textwidth]{plots/run9/unsafe_episode_rate}
        \caption{Unsafe episode rate.}
        \label{fig:run9_unsafe_rate}
    \end{subfigure}

    \vspace{0.5em}
    \begin{subfigure}[b]{0.48\textwidth}
        \centering
        \includegraphics[width=\textwidth]{plots/run9/intervention_rate}
        \caption{SafeDAgger intervention rate $\beta_{\text{safe}}$.}
        \label{fig:intervention_rate}
    \end{subfigure}
    \caption{Training curves comparing SafeDAgger (blue), DAgger (orange), and BC (green). (a)~Lift success relative to the teacher. (b)~Unsafe episode rate. (c)~SafeDAgger intervention rate, decaying from $\sim$0.7 to $\sim$0.1 within 20k iterations. DAgger has no intervention ($\beta = 0$); BC always uses the teacher ($\beta = 1$).}
    \label{fig:run9_training_curves}
\end{figure}

\subsubsection{Safety Analysis}%
\label{sec:safety_analysis}%

The three-way comparison reveals a clear trade-off between task performance and training-time safety.
BC maximizes lift success but produces the most unsafe student at evaluation time (41.8\,\% \ac{UER}), because the student never practices recovery from its own mistakes during training.
DAgger reduces the \ac{UER} to 32.7\,\% through online distribution correction but allows the student to visit dangerous states freely.
SafeDAgger further reduces the \ac{UER} to 30.2\,\% by intercepting dangerous actions during training, while matching DAgger's task performance.
As shown in \cref{fig:intervention_rate}, the SafeDAgger intervention rate $\beta_{\text{safe}}$ decays from $\sim$0.7 to $\sim$0.1 within 20k iterations, indicating the student acts autonomously in $\sim$90\,\% of environments for the majority of training.

\Cref{fig:failure_modes_grid} shows the per-failure-mode comparison for three representative objects.
On the basketball shoe, DAgger's harmful collision rate remains persistently high ($\sim$25\,\%) throughout training, while SafeDAgger reduces it to $\sim$5\,\% after 20k iterations.
On the teddy bear, DAgger's object out of bounds rate fluctuates between 30--45\,\%, whereas SafeDAgger stabilizes at $\sim$15\,\%.
The toy bagger shows a dramatic early collision spike for DAgger ($\sim$55\,\%) that SafeDAgger avoids entirely through teacher intervention.
BC shows low training-time failure rates because the teacher drives the environment, but the student inherits none of this robustness at evaluation time.
The complete per-object failure mode breakdowns for all methods and objects are provided in \cref{app:per_object_failure}.

\begin{figure}[htbp]
    \centering
    % Row 1: Basketball shoe
    \begin{subfigure}[b]{0.32\textwidth}
        \centering
        \includegraphics[width=\textwidth]{plots/run9/single_object_reason/single_basketball_shoe_harmful_collision}
    \end{subfigure}\hfill
    \begin{subfigure}[b]{0.32\textwidth}
        \centering
        \includegraphics[width=\textwidth]{plots/run9/single_object_reason/single_basketball_shoe_object_out_of_bound}
    \end{subfigure}\hfill
    \begin{subfigure}[b]{0.32\textwidth}
        \centering
        \includegraphics[width=\textwidth]{plots/run9/single_object_reason/single_basketball_shoe_palm_flipped}
    \end{subfigure}

    \vspace{0.3em}
    % Row 2: Teddy bear
    \begin{subfigure}[b]{0.32\textwidth}
        \centering
        \includegraphics[width=\textwidth]{plots/run9/single_object_reason/single_teddy_bear_harmful_collision}
    \end{subfigure}\hfill
    \begin{subfigure}[b]{0.32\textwidth}
        \centering
        \includegraphics[width=\textwidth]{plots/run9/single_object_reason/single_teddy_bear_object_out_of_bound}
    \end{subfigure}\hfill
    \begin{subfigure}[b]{0.32\textwidth}
        \centering
        \includegraphics[width=\textwidth]{plots/run9/single_object_reason/single_teddy_bear_palm_flipped}
    \end{subfigure}

    \vspace{0.3em}
    % Row 3: Toy bagger
    \begin{subfigure}[b]{0.32\textwidth}
        \centering
        \includegraphics[width=\textwidth]{plots/run9/single_object_reason/single_toy_bagger_harmful_collision}
    \end{subfigure}\hfill
    \begin{subfigure}[b]{0.32\textwidth}
        \centering
        \includegraphics[width=\textwidth]{plots/run9/single_object_reason/single_toy_bagger_object_out_of_bound}
    \end{subfigure}\hfill
    \begin{subfigure}[b]{0.32\textwidth}
        \centering
        \includegraphics[width=\textwidth]{plots/run9/single_object_reason/single_toy_bagger_palm_flipped}
    \end{subfigure}
    \caption{Per-object failure mode decomposition over distillation iterations for three representative objects (rows: basketball shoe, teddy bear, toy bagger; columns: harmful collision, object out of bounds, palm flip). SafeDAgger (blue) consistently reduces failure rates compared to DAgger (orange), with the largest margins on collision-prone objects.}
    \label{fig:failure_modes_grid}
\end{figure}

\subsubsection{Visual Representation Quality}%
\label{sec:aux_loss_analysis}%

The student network includes an auxiliary head that regresses the 3D object position from stereo images (\cref{eq:total_loss}).
Since the vision backbone and regression target are identical across all three methods, comparing the auxiliary loss isolates the effect of the training state distribution on visual feature learning.
\Cref{fig:aux_loss} shows that SafeDAgger and BC converge to nearly identical auxiliary loss values ($\sim$0.011 and $\sim$0.010, corresponding to $\sim$1\,cm object localization error), while DAgger's loss remains $\sim$1.7$\times$ higher ($\sim$0.018).
This gap arises because the DAgger student drives the environment itself, producing arm and hand configurations that differ from the teacher's trajectories and therefore generating camera viewpoints that are harder for the vision backbone to interpret.
Despite this gap, all three methods achieve object localization errors on the order of 1--2\,cm, indicating that the stereo vision backbone learns useful spatial representations regardless of the data-collection strategy.

\begin{figure}[htbp]
    \centering
    \includegraphics[width=0.7\textwidth]{plots/run9/aux_loss_object_pos}
    \caption{Auxiliary object-position regression loss over distillation iterations. SafeDAgger (blue) and BC (green) converge to similar values ($\sim$0.01), while DAgger (orange) remains $\sim$1.7$\times$ higher due to off-distribution viewpoints from student-driven rollouts.}
    \label{fig:aux_loss}
\end{figure}

\subsubsection{Per-Object Performance Breakdown}%
\label{sec:per_object_results}%

The 8 evaluation objects span a range of geometries, from compact shapes (closed fist, teddy bear) to thin or irregular objects (toy bagger, basketball shoe).
\Cref{fig:per_object_eval} shows the per-object lift success and unsafe episode rates for all three methods.
Performance varies substantially across objects: compact shapes achieve high lift success ($>$80\,\%) for all methods, while asymmetric objects show larger gaps.
SafeDAgger consistently exhibits lower unsafe rates across most objects compared to both DAgger and BC.

\begin{figure}[htbp]
    \centering
    \begin{subfigure}[b]{0.85\textwidth}
        \centering
        \includegraphics[width=\textwidth]{plots/run9/per_object_lift_eval}
        \caption{Per-object lift success.}
        \label{fig:per_object_lift_eval}
    \end{subfigure}

    \vspace{0.5em}
    \begin{subfigure}[b]{0.85\textwidth}
        \centering
        \includegraphics[width=\textwidth]{plots/run9/per_object_unsafe_eval}
        \caption{Per-object unsafe rate.}
        \label{fig:per_object_unsafe_eval}
    \end{subfigure}
    \caption{Per-object standalone evaluation of BC (green), DAgger (orange), and SafeDAgger (blue). (a)~Lift success varies by object geometry, with compact shapes performing best. BC achieves the highest lift but at the cost of safety. (b)~SafeDAgger exhibits lower unsafe episode rates across most objects.}
    \label{fig:per_object_eval}
\end{figure}

\Cref{fig:per_object_lift_h2h,fig:per_object_unsafe_h2h} provide a more detailed view by showing the per-object training curves over distillation iterations.
The performance gap between methods is not uniform: on objects such as the teddy bear and milk pot, SafeDAgger converges to noticeably higher lift success, while on others (e.g., turtle candle holder) all three methods perform similarly.
The per-object unsafe rate curves show that SafeDAgger achieves consistently lower collision rates on most objects, with the largest safety margins on objects where DAgger and BC produce aggressive approach behaviors.

\begin{figure}[htbp]
    \centering
    \includegraphics[width=\textwidth]{plots/run9/per_object_lifted_head_to_head}
    \caption{Per-object lift success over distillation iterations for SafeDAgger (blue), DAgger (orange), and BC (green). Each subplot shows one of the 8 evaluation objects.}
    \label{fig:per_object_lift_h2h}
\end{figure}

\begin{figure}[htbp]
    \centering
    \includegraphics[width=\textwidth]{plots/run9/per_object_unsafe_episode_head_to_head}
    \caption{Per-object unsafe episode rate over distillation iterations for SafeDAgger (blue), DAgger (orange), and BC (green). SafeDAgger consistently achieves lower unsafe rates on most objects.}
    \label{fig:per_object_unsafe_h2h}
\end{figure}

Additional experiments on environment scaling and ADR-randomised distillation are reported in \cref{app:scaling_adr}.

%% =========================================================================
\section{Summary of Findings}%
\label{sec:results_summary}%

The key experimental findings are:

\begin{enumerate}
    \item \textbf{Online distribution correction is essential for safe distillation.}  BC achieves the highest \ac{GSR} (87.8\,\%), matching the teacher, but produces the highest \ac{UER} (41.8\,\%).  The student learns to grasp well from the teacher's state distribution but cannot maintain safe behavior when it drifts, confirming the covariate-shift problem in privileged policy distillation.
    \item \textbf{SafeDAgger matches DAgger's task performance while improving safety.}  SafeDAgger (72.7\,\% \ac{GSR}) and DAgger (71.7\,\%) recover roughly 85\,\% of the teacher ceiling with comparable task performance.  However, SafeDAgger achieves a lower \ac{UER} (30.2\,\% vs.\ 32.7\,\%) at evaluation and maintains $\sim$15--20\,\% unsafe rate during training compared to $\sim$30--40\,\% for DAgger.
    \item \textbf{SafeDAgger's safety advantage is most pronounced on collision-prone objects.}  The per-object failure mode decomposition reveals that SafeDAgger reduces harmful collision rates by up to 20 percentage points on objects such as the basketball shoe and toy bagger, where DAgger's unrestricted exploration causes persistent contact with the table surface.
    \item \textbf{BC overfits with extended training.}  Lift success peaks at 95.6\,\% after 50k iterations and degrades to 85.6\,\% at 100k, while the \ac{UER} remains consistently high (35--49\,\%) regardless of training duration.
    \item \textbf{Safety threshold $\delta = 2.0$ provides the best balance.}  Calibrated from the L2 action-disagreement distribution, $\delta = 2.0$ produces $\sim$30\,\% initial intervention decaying to $\sim$5\,\% at convergence, balancing student autonomy with teacher guidance on challenging states.
\end{enumerate}

Additional implementation findings related to physics stability and teacher--student observation mapping are discussed in \cref{sec:impl_details}.
%
